% ==============================================================================
% CAPÍTULO 2: FUNDAMENTAÇÃO TEÓRICA E ESTADO DA ARTE
% ==============================================================================
\section{FUNDAMENTAÇÃO TEÓRICA E ESTADO DA ARTE}

O desenvolvimento de um ambiente computacional integrado para o gerenciamento térmico-energético de Centros de Dados Orbitais (\textit{Orbital Data Centers} -- ODCs) exige a convergência interdisciplinar de modelos físicos e técnicas avançadas de engenharia de controle e automação. Este capítulo apresenta as bases conceituais, formulações matemáticas analíticas e as tecnologias de ponta que sustentam a proposta da dissertação.

\subsection{Ambiente Espacial e Condições de Contorno de Centros de Dados Orbitais}

A transferência de infraestruturas massivas de processamento de inteligência artificial da superfície terrestre para plataformas orbitais fundamenta-se nas características singulares do ambiente exoatmosférico. Em solo, a radiação solar sofre expressiva atenuação, dispersão de Rayleigh e absorção por vapor de água, dióxido de carbono e ozônio, sendo tipicamente normalizada pelo padrão terrestre AM1.5 com irradiância média de pico de $1000\text{ W/m}^2$ (PIGNATON et al., 2021, p. 2). Ademais, a alternância entre dia e noite e a cobertura de nuvens impõem intermitência severa e baixo fator de capacidade a usinas fotovoltaicas em solo.

No espaço exterior, acima da atmosfera terrestre, a densidade de fluxo de radiação solar direta é governada pela constante solar extraterrestre sob o espectro AM0 (\textit{Air Mass Zero}), cujo valor médio nominal padronizado é dado por:
\begin{equation}
    G_{AM0} \approx 1361\text{ W/m}^2
\end{equation}
com variação anual inferior a $\pm 3{,}4\%$ decorrente da excentricidade da órbita elíptica da Terra em torno do Sol (WERTZ; EVERETT; PUSCHELL, 2011, p. 410).

Para viabilizar a operação contínua e economicamente competitiva de servidores de alta densidade sem a necessidade de armazenar dezenas de megawatts-hora em baterias pesadas, a literatura técnica e iniciativas da indústria aeroespacial apontam três opções orbitais clássicas (AGÜERA Y ARCAS et al., 2025, p. 2; BLOCH et al., 2026, p. 4):
\begin{enumerate}[label=(\alph*)]
    \item \textbf{Órbita Baixa Terrestre Convencional (\textit{Low Earth Orbit} -- LEO):} Para o cenário de estresse adotado nesta dissertação, considera-se órbita circular a $500\text{ km}$ e inclinação $i=51{,}6^\circ$. A duração do eclipse não é determinada pela altitude ou inclinação isoladamente: depende também da geometria Sol--plano orbital, expressa pelo ângulo beta $\beta$, e da época orbital. No caso de eclipse máximo empregado no Capítulo~4 ($\beta\approx0^\circ$), o modelo com sombra cônica calcula cerca de $35{,}8\text{ min}$ por revolução [CALCULADO]. Esse número é um resultado do cenário, não uma duração universal para toda órbita LEO. As temperaturas de $-150^\circ\text{C}$ a $+120^\circ\text{C}$ não são adotadas aqui como limites genéricos do satélite: a temperatura de cada componente deve resultar do modelo térmico, das propriedades ópticas, da atitude, das capacitâncias e das condições de contorno [PROVISÓRIO].
    \item \textbf{Órbita Heliossíncrona (\textit{Sun-Synchronous Orbit} -- SSO) do tipo \textit{Dawn-Dusk}:} O estudo adota como cenário nominal de referência uma órbita circular a $700\text{ km}$ e LTAN $18{:}00$. A faixa $85^\circ\leq |\beta|\leq90^\circ$ é tratada somente como valor nominal informado para uma condição orbital específica [PROVISÓRIO], não como envelope anual. Em órbitas Dawn-Dusk, a iluminação é quase contínua durante a maior parte do ano, mas podem ocorrer duas breves temporadas de eclipses sazonais associadas à obliquidade terrestre (SPIETZ et al., 2021). Assim, a avaliação anual requer calcular $\beta(t)$ com a época/efeméride e aplicar a Equação~(\ref{eq:eclipse_circular_beta}); a existência ou duração de eclipse não deve ser fixada apenas pela altitude.
    \item \textbf{Órbita Geoestacionária (\textit{Geostationary Earth Orbit} -- GEO):} Considera-se o caso comparativo equatorial a $35.786\text{ km}$. A iluminação é contínua fora das temporadas de eclipse; perto dos equinócios ocorrem eclipses diários, cuja duração máxima é da ordem de $72\text{ min}$ [VERIFICÁVEL] (HANLEY, 1979). A exposição aos cinturões de radiação e os requisitos de inserção dependem da trajetória de transferência e da missão; por isso, não se adota aqui um multiplicador universal de custo ou de $\Delta v$.
\end{enumerate}

Para uma primeira estimativa geométrica de eclipse em órbita circular, a dependência com o ângulo beta pode ser representada por (ESCOBAL; JOHNSON, 1972):
\begin{equation}
    T_{orb}=2\pi\sqrt{\frac{a^3}{\mu_E}}, \qquad
    t_{ecl}(\beta)=\frac{T_{orb}}{\pi}\cos^{-1}\!\left(\frac{\sqrt{1-(R_E/a)^2}}{\cos\beta}\right), \quad a=R_E+h,
    \label{eq:eclipse_circular_beta}
\end{equation}
quando $|\beta|<\sin^{-1}(R_E/a)$; fora dessa condição, o modelo de sombra cilíndrica prevê eclipse umbral nulo. Aqui, $R_E=6371\text{ km}$ e $\mu_E\approx398600\text{ km}^3/\text{s}^2$ (NASA, 2024). Para $h=500\text{ km}$ e $\beta=0^\circ$, essa aproximação fornece $T_{orb}\approx94{,}5\text{ min}$ e $t_{ecl}\approx35{,}7\text{ min}$ [CALCULADO]; o valor de $35{,}8\text{ min}$ citado acima decorre do modelo cônico com penumbra empregado no Capítulo~4. O modelo geométrico não substitui a propagação orbital com efemérides para uma missão específica.

Neste trabalho, a SSO Dawn-Dusk é o \textbf{cenário nominal de referência}, enquanto LEO convencional e GEO permanecem como cenários comparativos independentes e de estresse. A redução do eclipse na referência SSO pode reduzir a energia de armazenamento requerida, mas não se afirma uma redução percentual universal da massa de baterias: esse resultado depende do perfil de carga, da potência fotovoltaica no fim de vida, do ângulo beta, da eficiência de conversão e dos limites de SoC/DoD, e deve ser calculado para cada cenário.

Outro aspecto estrutural determinante consiste na arquitetura da plataforma espacial. Pela relação de massa conceitual desta seção, cada nó adotado de $50\text{ kW}$ tem massa estimada de $2.600\text{ kg}$; os 22 nós fabricados e lançados (20 operacionais e 2 reservas) totalizam $57{,}2\text{ t}$ [CALCULADO], sem adicionar massas de adaptadores ou margens não representadas no modelo. Portanto, a estimativa de aproximadamente $100\text{ t}$ não é compatível com a relação de massa usada nesta tabela e não é adotada como resultado do projeto. O estudo mantém a arquitetura de \textbf{clusters modulares distribuídos} como premissa conceitual; não afirma que um lançador específico seja incapaz de transportar uma massa determinada para SSO sem análise de desempenho da missão. Nesses clusters, os módulos de processamento cooperam através de \textbf{Enlaces Ópticos Inter-Satélite (\textit{Optical Inter-Satellite Links} -- OISL)}, cujos transceivers a laser estabelecem redes de interconexão de alta velocidade (100~Gbps a 1~Tbps) para sincronização de gradientes e comunicação inter-nós (AGÜERA Y ARCAS et al., 2025, p. 5; BLOCH et al., 2026, p. 7; KAUSHAL; KADDOUM, 2017, p. 58).

\subsubsection{Arquitetura de Cluster e Trade-Off Técnico-Econômico de Escala}
\label{sec:tradeoff_custo_escala}

A determinação da potência unitária de cada nó satelital ($P_{no}$) é tratada como um trade-off paramétrico para atender a uma capacidade operacional de referência de pelo menos $1{,}0\text{ MW}$ de geração bruta. Na arquitetura adotada, 20 nós operacionais de $50\text{ kW}$ fornecem $1{,}0\text{ MW}$ bruto; cada nó destina $40\text{ kW}$ à carga de TI e os $10\text{ kW}$ restantes à recarga e às cargas auxiliares, totalizando $800\text{ kW}$ úteis de TI. Essa relação de potência útil é específica da arquitetura de $50\text{ kW}$ e não é extrapolada para os demais tamanhos avaliados. Com redundância de $10\%$, a arquitetura adotada lança 22 nós, correspondendo a $1{,}1\text{ MW}$ de capacidade bruta nominal instalada, dos quais dois nós são reserva.

Para cada alternativa, o número de nós operacionais requerido e o total fabricado/lançado são:
\begin{equation}
    N_{req}=\left\lceil\frac{1000}{P_{no}}\right\rceil, \qquad
    N_{tot}=\left\lceil\frac{11}{10}N_{req}\right\rceil
    =\left\lfloor\frac{11N_{req}+9}{10}\right\rfloor, \qquad
    P_{bruta,op}=N_{req}P_{no},
    \label{eq:nos_tradeoff}
\end{equation}
onde $P_{no}$ é expresso em kW, $P_{bruta,op}$ é a potência bruta dos nós operacionais e $N_{tot}$ inclui a margem de redundância. A forma inteira equivalente de $N_{tot}$ evita erros de arredondamento numérico em implementações computacionais.

A comparação econômica limita-se à \textbf{Estimativa Preliminar de Custo de Implantação Espacial (CAPEX Orbital: Fabricação + Lançamento)}:
\begin{equation}
    C_{CAPEX}(P_{no})=C_{mfg}(P_{no},N_{tot})+C_{launch}(P_{no},N_{tot}).
    \label{eq:capex_orbital}
\end{equation}
Os custos operacionais contínuos de estação de solo, manutenção de constelação e reposição em órbita ($C_{ops}$) não foram quantificados; portanto, não são incluídos no CAPEX apresentado e ficam como trabalho futuro. A estimativa tampouco inclui custos não recorrentes de desenvolvimento, qualificação, adaptadores especiais ou análise de integração de lançamento.

A massa de cada veículo é aproximada pela relação paramétrica $M_{sat}=M_0+\beta_mP_{no}$, com $M_0\approx350\text{ kg}$ e $\beta_m\approx45\text{ kg/kW}$ [PROVISÓRIO]. Esses coeficientes representam uma hipótese de arquitetura conceitual e não uma massa de satélite qualificada. O custo de fabricação do lote é estimado por:
\begin{equation}
    C_{mfg}=\left(C_{fixo}+\gamma_pP_{no}\right)N_{tot}^{1+b},
    \label{eq:custo_fabricacao}
\end{equation}
com $C_{fixo}=\text{US\$ }5{,}0\text{ M}$, $\gamma_p=\text{US\$ }90\text{ k/kW}$, taxa de progresso $S=85\%$ e $b=\log_2(S)\approx-0{,}234$ [PROVISÓRIO]. A curva de aprendizagem é usada como modelo paramétrico simplificado. O próprio NASA Cost Estimating Handbook ressalta que a aplicação de curvas de aprendizagem ao setor espacial pode ser questionável quando se fabricam poucos itens (NASA, 2015, apêndice Q, p. Q-5); por isso, $S=85\%$ não é tratado como parâmetro validado para esta constelação. Os valores específicos de $C_{fixo}$, $\gamma_p$ e $S$ são premissas do cenário do autor, não coeficientes extraídos diretamente de NASA (2015) ou de WERTZ; EVERETT; PUSCHELL (2011). Os custos calculados são resultados dessa hipótese e não cotações de fabricação.

Para o lançamento, mantém-se $P_{voo}=\text{US\$ }67\text{ M}$ como premissa de cenário de mercado informada para o período 2021--2024 [PROVISÓRIO], com $C_{launch}=P_{voo}\left\lceil N_{tot}/n_{coifa}\right\rceil$. O guia técnico SpaceX de 2021 citado anteriormente descreve interfaces e requisitos do veículo, mas não constitui, por si só, fonte pública de preço. A ficha pública da SpaceX consultada apresenta US\$~69{,}75~M como preço-padrão para contratos até 2024; assim, o valor de US\$~67~M é mantido apenas como cenário de mercado do autor, não como preço atual ou preço oficial verificado (SPACEX, 2024). A aplicação do valor por voo dedicado é uma simplificação do modelo econômico.

A quantidade $n_{coifa}$ de satélites por voo e os limites de massa da Tabela~\ref{tab:tradeoff_custo_constelacao} são \textbf{hipóteses conceituais de empacotamento} [PROVISÓRIO]. A coifa Falcon tem diâmetro externo de $5{,}2\text{ m}$, mas a compatibilidade de cada carga requer verificar o volume estático de acomodação, adaptadores, centro de gravidade, deflexões dinâmicas e análise acoplada de cargas. O guia da SpaceX apresenta limites dependentes da configuração e determina coordenação/análise específica para missões com múltiplas cargas (SPACEX, 2025). Não foram realizados CAD/FEA nem ensaios de integração/separação para os satélites desta arquitetura. Portanto, os valores de satélites por coifa e o degrau de empacotamento a partir de $70\text{ kW}$ são premissas do trade-off, não demonstração de viabilidade mecânica. Também não se afirma capacidade de carga específica do Falcon~9 para a órbita SSO de $700\text{ km}$ sem análise de desempenho da missão.

Na Tabela~\ref{tab:tradeoff_custo_constelacao}, a meta de $1{,}0\text{ MW}$ brutos refere-se à potência nominal dos nós em operação. A capacidade nominal total da frota, $P_{bruta,tot}=N_{tot}P_{no}$, inclui a redundância; por isso, é apresentada separadamente de $P_{bruta,op}=N_{req}P_{no}$. O custo por kW bruto é calculado sobre $P_{bruta,op}$ para comparar alternativas que atendem à mesma capacidade operacional de referência, e não sobre potência útil de TI.

\begin{table}[htbp]
\centering
\caption{Trade-off para capacidade agregada de $1{,}0\text{ MW}$ brutos instalados em operação; a capacidade nominal total da frota inclui 10\% de redundância.}
\label{tab:tradeoff_custo_constelacao}
\resizebox{\textwidth}{!}{%
\begin{tabular}{cccccccccccc}
\hline
\textbf{$P_{no}$} & \textbf{$N_{req}$} & \textbf{$N_{tot}$} & \textbf{$P_{bruta,op}$ (kW)} & \textbf{$P_{bruta,tot}$ (kW)} & \textbf{Massa/sat. est.} & \textbf{Sats./coifa} & \textbf{Voos} & \textbf{$C_{mfg}$ (US\$ M)} & \textbf{$C_{launch}$ (US\$ M)} & \textbf{CAPEX orbital (US\$ M)} & \textbf{Custo / kW Bruto (US\$ M/kW)} \\ \hline
10 kW & 100 & 110 & 1.000 & 1.100 & 800 kg & 8 & 14 & 215,6 & 938,0 & 1.153,6 & 1,15 \\
20 kW & 50 & 55 & 1.000 & 1.100 & 1.250 kg & 5 & 11 & 146,2 & 737,0 & 883,2 & 0,88 \\
25 kW & 40 & 44 & 1.000 & 1.100 & 1.475 kg & 5 & 9 & 131,4 & 603,0 & 734,4 & 0,73 \\
35 kW & 29 & 32 & 1.015 & 1.120 & 1.925 kg & 3 & 11 & 115,7 & 737,0 & 852,7 & 0,84 \\
\textbf{50 kW (Adotado)} & \textbf{20} & \textbf{22} & \textbf{1.000} & \textbf{1.100} & \textbf{2.600 kg} & \textbf{2} & \textbf{11} & \textbf{101,2} & \textbf{737,0} & \textbf{838,2} & \textbf{0,84} \\
60 kW & 17 & 19 & 1.020 & 1.140 & 3.050 kg & 2 & 10 & 99,1 & 670,0 & 769,1 & 0,75 \\
70 kW & 15 & 17 & 1.050 & 1.190 & 3.500 kg & 1 & 17 & 98,9 & 1.139,0 & 1.237,9 & 1,18 \\
100 kW & 10 & 11 & 1.000 & 1.100 & 4.850 kg & 1 & 11 & 87,8 & 737,0 & 824,8 & 0,82 \\ \hline
\end{tabular}%
}
\end{table}

Na linha de $20\text{ kW}$, a contagem corrigida corresponde a 55 satélites e 220 terminais OISL, mantendo-se a fração de massa inerte estimada em $350/1250=28\%$ sob o modelo linear adotado. O menor CAPEX numérico da tabela ocorre em $60\text{ kW}$ (US\$~769{,}1~M), condicionado às premissas de custo e empacotamento; a transição para $70\text{ kW}$ eleva o CAPEX estimado para US\$~1.237{,}9~M (+61{,}0\%), principalmente porque a hipótese da tabela reduz de dois para um satélite por voo. Essa diferença é uma consequência do degrau discreto assumido, não uma conclusão validada por integração mecânica.

Apesar do menor CAPEX paramétrico em $60\text{ kW}$, adota-se $50\text{ kW}$ como ponto de projeto conceitual para os modelos eletrotérmicos. No modelo do Capítulo~4, o nó de $50\text{ kW}$ requer área radiativa efetiva ativa de $88{,}6\text{ m}^2$ sob as condições de referência. Se a carga térmica crescer linearmente com $P_{no}$ e forem mantidas as mesmas propriedades ópticas, temperatura e capacidade líquida de rejeição, a extrapolação para $60\text{ kW}$ resulta em $A_{rad,eff,60}=88{,}6(60/50)\approx106{,}3\text{ m}^2$ [CALCULADO; extrapolação de primeira ordem]. Esse valor indica um possível agravamento do empacotamento, mas não prova que o conjunto excede a coifa: área radiativa efetiva não é, por si só, área projetada ou envelope dobrado, e a verificação exige geometria implantada, análise de sombras/interferências e integração mecânica.

Como segundo critério de seleção, considera-se [PROVISÓRIO] uma arquitetura de conversores GaN multifásicos de $3\times20\text{ kW}$. Uma carga de $40\text{ kW}$ de TI corresponde a $66{,}7\%$ da potência nominal agregada de $60\text{ kW}$; porém, se os $10\text{ kW}$ de recarga/auxiliares também atravessarem esses conversores, a carga total de $50\text{ kW}$ corresponde a $83{,}3\%$. Com a perda de um módulo de $20\text{ kW}$, os dois módulos restantes fornecem no máximo $40\text{ kW}$ nominais; a manutenção da carga de TI exigiria descarte ou alimentação independente das cargas auxiliares/recarga, sem margem nominal para sobrecarga. A topologia, os limites térmicos e a estratégia de contingência dependem de seleção de componentes e validação posteriores.

Assim, $50\text{ kW}$ é o \textbf{ponto de projeto adotado} por compatibilidade com o modelo térmico e pela modularidade conceitual da conversão de potência; $60\text{ kW}$ é o menor resultado de CAPEX sob as hipóteses atuais. A viabilidade mecânica do lançamento duplo em $50\text{ kW}$ e a inviabilidade em $60\text{ kW}$ permanecem hipóteses [PROVISÓRIO] até análise CAD/FEA/CLA e definição do adaptador. Os valores da tabela são [CALCULADO] a partir de coeficientes e preços [PROVISÓRIO]; nenhum deles representa validação experimental ou cotação vinculante.

\subsection{O Gargalo da Transferência Térmica no Vácuo e a Dinâmica de Chips de IA}

Conquanto a disponibilidade solar seja abundante em órbitas heliossíncronas ou geoestacionárias, a operação contínua de semicondutores de alto desempenho esbarra em um limite termodinâmico intransponível: a transferência de calor em confinamento a vácuo (MESEGUER; PÉREZ-GRANDE; SANZ-ANDRÉS, 2012, p. 45).

Em data centers convencionais em terra, o calor gerado pelos circuitos integrados é extraído por condução até trocadores de calor e transferido para o ar ou água por convecção forçada monofásica ou bifásica (BARROSO; HÖLZLE; RANGANATHAN, 2019, p. 68). No vácuo orbital (pressão residual $P < 10^{-6}\text{ Pa}$), inexiste qualquer meio fluido contínuo para propiciar convecção ou condução externa. Consequentemente, \textbf{todo o calor gerado internamente pela carga útil computacional deve ser conduzido estruturalmente até a carcaça da plataforma e rejeitado para o espaço profundo estritamente por meio de radiação térmica superficial} (CHEN et al., 2026, p. 2).

A taxa de rejeição radiativa é matematicamente descrita pela Lei de Stefan-Boltzmann:
\begin{equation}
    q_{rad}(t) = \varepsilon \sigma A_{rad} \left[ T_{rad}(t)^4 - T_{space}^4 \right]
    \label{eq:stefan_boltzmann}
\end{equation}
em que:
\begin{itemize}
    \item $\varepsilon$ é a emissividade no infravermelho da cobertura óptica do radiador ($0 < \varepsilon \le 1$, tipicamente $\approx 0{,}85$ a $0{,}92$ para tintas térmicas brancas e filmes OSR -- \textit{Optical Solar Reflectors});
    \item $\sigma = 5{,}670 \times 10^{-8}\text{ W/m}^2\text{K}^4$ é a constante universal de Stefan-Boltzmann;
    \item $A_{rad}$ representa a área efetiva superficial de rejeição radiativa orientada para o espaço profundo ($\text{m}^2$);
    \item $T_{rad}(t)$ é a temperatura absoluta da superfície externa do radiador em Kelvin (K);
    \item $T_{space} \approx 3\text{ K}$ é a temperatura efetiva da radiação cósmica de fundo do espaço profundo.
\end{itemize}

O circuito elétrico equivalente de parâmetros concentrados (\textit{lumped thermal model}) que modela o nó de silício do acelerador de IA (GPU/TPU) e o radiador é expresso pelo sistema de equações diferenciais acopladas:
\begin{align}
    C_{j} \frac{dT_{j}(t)}{dt} &= P_{chip}(t) - \frac{T_{j}(t) - T_{rad}(t)}{R_{th,j-rad}} \label{eq:nodal_chip}\\
    C_{rad} \frac{dT_{rad}(t)}{dt} &= \frac{T_{j}(t) - T_{rad}(t)}{R_{th,j-rad}} - \varepsilon \sigma A_{rad} \left[ T_{rad}(t)^4 - T_{space}^4 \right] \label{eq:nodal_rad}
\end{align}
onde $T_j(t)$ é a temperatura da junção do semicondutor, $C_j$ e $C_{rad}$ são as capacitâncias térmicas do chip e da placa do radiador ($\text{J/K}$), $R_{th,j-rad}$ é a resistência térmica condutiva interfacial ($\text{K/W}$), e $P_{chip}(t)$ representa a potência elétrica instantânea dissipada pelo processador durante o cálculo das redes neurais.

A Equação~(\ref{eq:stefan_boltzmann}) expõe a severidade do gargalo: para que a junção de silício mantenha integridade dielétrica e física, a temperatura máxima é limitada a $T_j \le 85^\circ\text{C}$ ($358{,}15\text{ K}$). Admitindo uma queda condutiva de $15^\circ\text{C}$ na interface interna, a temperatura máxima do radiador não excede $T_{rad} \approx 70^\circ\text{C}$ ($343{,}15\text{ K}$). Aplicando a Lei de Stefan-Boltzmann sob essas condições ideais:
\begin{equation}
    \frac{q_{rad}}{A_{rad}} \approx 0{,}90 \cdot \left(5{,}670 \times 10^{-8}\right) \cdot \left( 343{,}15^4 - 3^4 \right) \approx 708\text{ W/m}^2
\end{equation}

Assim, cada quilowatt elétrico de servidores de IA operando em regime permanente exige uma área radiativa desobstruída superior a $1{,}41\text{ m}^2$. Em um rack orbital consumindo de 5~kW a 20~kW, radiadores fixos demandariam dezenas de metros quadrados de painéis desdobráveis de altíssimo custo e peso de lançamento (CHEN et al., 2026, p. 3). Se as cargas computacionais forem despachadas sem sensibilidade térmica, o calor acumulado satura rapidamente a capacitância térmica $C_j$, deflagrando estrangulamento térmico por corte abrupto (\textit{thermal shutdown}) ou colapso irreversível dos chips de IA.

\subsection{Subsistema Fotovoltaico Espacial e Conversão Estática de Potência}

A geração primária em satélites é assegurada pelo subsistema elétrico (\textit{Electrical Power System} -- EPS), composto pelo arranjo fotovoltaico, conversores chaveados de condicionamento e distribuição e baterias de sustentação (AL HOSANI et al., 2021, p. 9125; BARBI, 2006, p. 18).

Para células fotovoltaicas espaciais de alta eficiência (células de tripla junção InGaP/InGaAs/Ge com rendimento $\eta > 30\%$), a relação corrente-tensão sob o espectro AM0 é modelada pelo circuito equivalente de um diodo e cinco parâmetros (PIGNATON et al., 2021, p. 2):
\begin{equation}
    I = I_{ph} - I_0 \left[ \exp\left( \frac{q (V + I R_s)}{a k_B T_{cel}} \right) - 1 \right] - \frac{V + I R_s}{R_{sh}}
    \label{eq:pv_1diodo}
\end{equation}
em que $I_{ph}$ é a fotocorrente gerada pela incidência solar AM0, $I_0$ é a corrente de saturação reversa do diodo, $q = 1{,}602 \times 10^{-19}\text{ C}$ é a carga elementar, $k_B = 1{,}381 \times 10^{-23}\text{ J/K}$ é a constante de Boltzmann, $a$ é o fator de idealidade da junção, $T_{cel}$ é a temperatura das células em Kelvin, e $R_s$ e $R_{sh}$ representam as resistências série e paralelo (shunt), respectivamente.

A dependência térmica e de irradiância é expressa por:
\begin{align}
    I_{ph} &= \left( I_{ph,ref} + K_I \Delta T \right) \frac{G_{AM0}}{G_{ref}} \\
    I_0 &= I_{0,ref} \left( \frac{T_{cel}}{T_{ref}} \right)^3 \exp\left[ \frac{q E_g}{a k_B} \left( \frac{1}{T_{ref}} - \frac{1}{T_{cel}} \right) \right]
\end{align}
sendo $K_I$ o coeficiente de temperatura da corrente e $E_g$ o \textit{bandgap} do semicondutor.

Para garantir que a matriz solar opere continuamente em seu Ponto de Máxima Potência (MPP) frente às flutuações de temperatura causadas por manobras de atitude e variações transitórias, conversores DC-DC integrados executam algoritmos avançados de rastreamento (\textit{Maximum Power Point Tracking} -- MPPT), incluindo arquiteturas baseadas em aprendizado profundo e técnicas adaptativas (ABAGERO et al., 2025, p. 2; SHOBHA; USHA; KUMAR, 2023, p. 69). 

No estágio de distribuição de potência do satélite, o barramento primário de alta tensão em corrente contínua (HVDC de $400\text{ VDC}$), dimensionado para minimizar drasticamente as perdas ôhmicas de condução e a massa de condutores sem atingir os limiares de descarga parcial e efeito corona no vácuo espacial estipulados pela norma NASA-HDBK-4007A (NASA, 2026, p. 40), deve ser rebaixado para os níveis intermediários de distribuição de rack ($48\text{ VDC}$) e subsequentemente para as tensões sub-1~V demandadas pelos núcleos de processamento paralelo dos aceleradores de inteligência artificial ($0{,}8\text{ VDC}$). Essa conversão em múltiplos estágios é executada por conversores DC-DC modulares de alta densidade baseados em semicondutores de nitreto de gálio (GaN) com retificação síncrona (BARBI, 2006, p. 65; SARFI; BEIK, 2024).

As perdas de potência internas dissipadas nesses conversores são descritas analiticamente por:
\begin{equation}
    P_{loss,conv} = P_{cond} + P_{sw} + P_{core} = I_{rms}^2 R_{on} + \frac{1}{2} V_{in} I_{out} (t_r + t_f) f_{sw} + P_{core}
    \label{eq:perdas_conversores}
\end{equation}
onde $R_{on}$ é a resistência de condução dos transistores GaN/MOSFET, $t_r$ e $t_f$ são os tempos de subida e descida de chaveamento, $f_{sw}$ é a frequência de comutação e $P_{core}$ são as perdas magnéticas nos indutores de filtro. Como esses conversores operam confinados na carcaça do rack, $P_{loss,conv}$ atua como uma fonte secundária de aquecimento que onera diretamente o balanço térmico interno, exigindo eficiências superiores a 95\% ($\eta_{conv} > 0{,}95$).

Adicionalmente, em uma arquitetura modular de múltiplos satélites cooperantes em cluster, o subsistema de potência elétrica deve suprir a infraestrutura de intercomunicação óptica inter-satélite (\textit{Optical Inter-Satellite Links} -- OISL). Cada terminal laser consome tipicamente entre $60\text{ W}$ e $150\text{ W}$ de potência elétrica contínua para alimentação dos diodos de bombeamento e dos espelhos de direcionamento fino (\textit{Fast Steering Mirrors} -- FSM). Com rendimento eletro-óptico situado entre $20\%$ e $30\%$ ($\eta_{eo} \approx 0{,}20$ a $0{,}30$), cerca de $70\%$ a $80\%$ da energia consumida pelos módulos OISL é convertida em calor residual interno ($P_{th,OISL} = (1 - \eta_{eo}) P_{OISL}$), que se soma às perdas dos conversores e dos chips de IA. Assim, o balanço de potência do barramento DC é formulado como:
\begin{equation}
    P_{total}(t) = P_{IA}(t) + P_{OISL}(t) + P_{aux}(t) + P_{loss,conv}(t)
    \label{eq:balanco_potencia_total}
\end{equation}
onde $P_{aux}$ agrega o consumo dos computadores de bordo (OBC), telemetria e bombas térmicas.

\subsection{Tecnologia Central: Sistema de Gerenciamento Térmico-Energético (TEMS)}

A tecnologia central investigada no nível de mestrado consiste no \textbf{Sistema de Gerenciamento Térmico-Energético (\textit{Thermal-Energy Management System} -- TEMS)}. Diferentemente de sistemas terrestres — onde a climatização independe da lógica de software e opera de maneira desacoplada via chillers de expansão —, no vácuo espacial o software de processamento e a rejeição física de calor formam um sistema dinâmico intrinsecamente acoplado (CHEN et al., 2026, p. 4; SINGH, 2026, p. 2).

O TEMS orquestra duas variáveis de controle primárias em tempo real:
\begin{enumerate}
    \item \textbf{Modulação Dinâmica de Tensão e Frequência (\textit{Dynamic Voltage and Frequency Scaling} -- DVFS):} Ajuste da frequência de relógio $f_{clk}$ e da tensão de núcleo $V_{core}$ de cada GPU/TPU de IA. A potência dinâmica dissipada pelo chip segue a relação cúbica:
    \begin{equation}
        P_{dyn} = \alpha C_{eff} V_{core}^2 f_{clk}
    \end{equation}
    onde $\alpha$ é o fator de atividade e $C_{eff}$ é a capacitância efetiva de chaveamento lógico. A redução preventiva de $f_{clk}$ mitiga quadraticamente a taxa de injeção de calor $P_{chip}$ no nó térmico da Equação~(\ref{eq:nodal_chip}).
    \item \textbf{Escalonamento Térmico-Consciente de Tarefas de IA (\textit{Thermal-Aware Workload Scheduling}):} Alocação e sequenciamento de lotes de treinamento e inferência de redes neurais profundas entre diferentes nós computacionais do rack orbital.
\end{enumerate}

O problema de controle ótimo formulado no TEMS é expresso como a minimização da função custo multiobjetivo em horizonte deslizante:
\begin{equation}
\begin{split}
    \min_{\mathbf{u}[k]} J = \sum_{k=0}^{N-1} \Big\{ & w_1 \left[ T_{j}[k] - T_{ref} \right]^2 + w_2 \Phi_{penalidade}(T_{j}[k] - T_{max}) \\
    & + w_3 P_{loss,conv}[k] - w_4 \Psi_{throughput}[k] \Big\}
\end{split}
\label{eq:funcao_custo_tems}
\end{equation}
sujeito às restrições de estado e controle:
\begin{align}
    T_{j}[k] &\le T_{max} = 85^\circ\text{C}, \quad \forall k \\
    P_{total}[k] &\le P_{mpp}[k] + P_{bat,max}[k] \\
    \mathbf{u}[k] &\in \left\{ f_{min} \le f_{clk} \le f_{max}, \, \text{Modo}_{k} \in \{\text{Idle}, \text{Inferência}, \text{Treinamento}, \text{Downlink}\} \right\}
\end{align}
sendo $w_1, w_2, w_3, w_4$ fatores de ponderação, $\Phi_{penalidade}$ uma função barreira para violação de limites térmicos e $\Psi_{throughput}$ a vazão líquida de operações de ponto flutuante (FLOPs) executadas por unidade de tempo.

O algoritmo do TEMS opera sob uma Máquina de Estados Finitos com transições supervisionadas, garantindo que nenhum nó atinja saturação térmica enquanto mantém a máxima eficiência de uso de energia equivalente orbital ($PUE_{orbital} \rightarrow 1{,}05$). Essa formulação apoia-se e expande o arcabouço de Co-Design 4D em malha fechada introduzido pelo autor em Barros e Fernandes (2026a), no qual o controle preditivo econômico não-linear (4D-NMPC) de ordem reduzida (ROM) acopla estritamente o balanço térmico de semicondutores, a estabilização do barramento elétrico, o controle de atitude (ADCS) e a vazão de dados com compressão semântica em banda Ka.

\subsection{Tecnologia Complementar: Estimativa de Estado de Carga (SoC) via Filtro de Kalman Adaptativo}

Embora a órbita heliossíncrona SSO \textit{dawn-dusk} proporcione insolação direta contínua, bancos de baterias eletroquímicas de íons de lítio (ou polímero de lítio) são componentes indispensáveis para: (i) sustentação do barramento durante manobras orbitais em que os painéis perdem temporariamente a visada solar; (ii) filtragem de picos transitórios de corrente ($\text{kA}/\mu\text{s}$) exigidos por disparos de nós de processamento neural; e (iii) redundância vital de segurança operacional (NASA, 2026, p. 40; CAPPELLETTI; BATTISTINI; MALPHRUS, 2020, p. 112).

\subsubsection{Envelope Térmico e Limites Físico-Químicos do Banco de Baterias Espacial}

Diferentemente da carcaça e dos painéis solares externos — que no vácuo oscilam amplamente entre $-150^\circ\text{C}$ e $+120^\circ\text{C}$ —, as células de bateria de lítio operam estritamente confinadas no interior do chassi, protegidas por mantas isolantes de múltiplas camadas (\textit{Multi-Layer Insulation} -- MLI) e aquecedores elétricos resistivos de sobrevivência (\textit{patch heaters}) regulados por termostato (conforme as diretrizes da NASA SP-8005 e ECSS-E-ST-20-40C). Esse isolamento garante que o banco opere em uma faixa nominal estreita de $+10^\circ\text{C}$ a $+25^\circ\text{C}$, com limites extremos de contingência e qualificação de $-20^\circ\text{C}$ a $+50^\circ\text{C}$.

Essa restrição térmica decorre de limitações físico-químicas severas:
\begin{enumerate}
    \item \textbf{Carga sob Temperaturas Negativas ($T < 0^\circ\text{C}$):} A cinética de difusão do lítio nos sítios intersticiais do ânodo de grafite torna-se excessivamente lenta frente à corrente aplicada, deflagrando o fenômeno de deposição metálica de lítio (\textit{lithium plating}). Essa deposição forma dendritos que perfuram o separador polimérico, provocando curto-circuito interno irreversível e falha catastrófica da missão.
    \item \textbf{Operação em Temperaturas Elevadas ($T > 45^\circ\text{C}$):} Ocorre a degradação acelerada do eletrólito líquido, espessamento da camada de passivação de interface (\textit{Solid Electrolyte Interphase} -- SEI) e perda de lítio reciclável, podendo culminar em escape térmico (\textit{thermal runaway}).
\end{enumerate}

Mesmo dentro do envelope de controle interno ($+5^\circ\text{C}$ a $+35^\circ\text{C}$), a resistência ôhmica interna $R_0(T)$ exibe forte dependência de temperatura, governada pela cinética da Equação de Arrhenius:
\begin{equation}
    R_0(T) = R_{0,ref} \exp\left[ \frac{E_a}{k_B} \left( \frac{1}{T_{bat}} - \frac{1}{T_{ref}} \right) \right]
\end{equation}
onde $E_a$ é a energia de ativação eletroquímica. Uma redução térmica de $+25^\circ\text{C}$ para $0^\circ\text{C}$ triplica a resistência interna, alterando drasticamente a curva terminal $V_t$ e induzindo erros severos em modelos de parâmetros fixos.

\subsubsection{Justificativa do Algoritmo AUKF frente ao Estado da Arte}

A literatura científica recente apresenta quatro grandes famílias para estimativa de Estado de Carga (SoC):
\begin{enumerate}
    \item \textbf{Métodos Empíricos (Contagem de Coulomb e Tabela OCV):} Sofrem de desvio cumulativo (\textit{drift}) por ruído de sensores de corrente e exigem longos períodos de repouso (1~h a 4~h) para relaxação da tensão de circuito aberto, inviáveis em servidores de operação ininterrupta.
    \item \textbf{Modelos Eletroquímicos de Primeiros Princípios (P2D e SPM):} Modelam as equações diferenciais parciais de difusão de Fick e cinética de Butler-Volmer. Conquanto apresentem máxima fidelidade físico-química, sua discretização espacial demanda centenas de variáveis de estado, tornando a execução em tempo real inviável para processadores espaciais de baixo consumo tolerantes a radiação (e.g., arquiteturas RAD750 a 133~MHz).
    \item \textbf{Redes Neurais e Aprendizado Profundo Puro (LSTM, GRU, Transformers):} Conquanto alcancem elevada precisão em ensaios terrestres padronizados, comportam-se como caixas-pretas estatísticas. Diante de condições operacionais orbitais inéditas (degradação por radiação cósmica pesada ou variações térmicas anômalas), tais modelos podem divergir bruscamente ou prever valores espúrios ($SoC > 100\%$), colocando em risco o suprimento de energia do satélite.
    \item \textbf{Filtros Gaussianos Não Lineares Adaptativos (AUKF):} O Filtro de Kalman Unscented propaga $2L+1$ pontos sigma determinísticos através da não linearidade da curva $V_{oc}(SoC)$, capturando a média e a covariância com precisão de 3ª ordem sem requerer cálculo de derivadas Jacobianas (ao contrário do EKF, que trunca a série de Taylor em 1ª ordem e diverge sob degraus de corrente $\text{kA}/\mu\text{s}$). A incorporação de fator de esquecimento adaptativo $\lambda_k \ge 1$ (AUKF) confere prioridade às medições recentes e recalibra online as matrizes de covariância frente à deriva de $R_0(T)$ (FENG et al., 2022, p. 44522).
\end{enumerate}

Dessa forma, o **AUKF com modelo de circuito equivalente de Thévenin de 2ª ordem térmico-dependente (Dual Polarization)** estabelece o equilíbrio ótimo entre precisão analítica ($RMSE < 1{,}5\%$), estabilidade assintótica garantida pela Teoria de Lyapunov e tempo de execução na escala de microssegundos em hardware de bordo espacial.

O modelo elétrico equivalente de 2ª ordem adotado é expresso por:
\begin{align}
    V_t(t) &= V_{oc}(SoC) - I_L(t) R_0(T) - V_{p1}(t) - V_{p2}(t) \\
    \frac{dV_{p1}(t)}{dt} &= -\frac{1}{R_{p1} C_{p1}} V_{p1}(t) + \frac{1}{C_{p1}} I_L(t) \\
    \frac{dV_{p2}(t)}{dt} &= -\frac{1}{R_{p2} C_{p2}} V_{p2}(t) + \frac{1}{C_{p2}} I_L(t)
\end{align}
onde os pares $R_{p1}C_{p1}$ e $R_{p2}C_{p2}$ capturam a dinâmica de transferência de carga e a relaxação lenta de difusão de estado sólido. A equação de propagação com fator de esquecimento $\lambda_k \ge 1$ segue:
\begin{equation}
    \mathbf{P}_{k|k-1} = \lambda_k \sum_{i=0}^{2L} W_i^{(c)} \left[ \boldsymbol{\chi}_{i, k|k-1} - \hat{\mathbf{x}}_{k|k-1} \right] \left[ \boldsymbol{\chi}_{i, k|k-1} - \hat{\mathbf{x}}_{k|k-1} \right]^T + \mathbf{Q}_k
\end{equation}
Esse sinal de SoC alimenta continuamente o escalonador do TEMS para autorizar ou adiar ciclos pesados de treinamento neural.

\subsection{Síntese do Estado da Arte e Análise Comparativa da Literatura}

A consolidação da fundamentação teórica revela que a literatura científica recente aborda as dimensões do problema de forma fragmentada. O Quadro~\ref{tab:estado_arte} sintetiza o estado da arte e evidencia a lacuna precisa que a presente dissertação se propõe a preencher.

\begin{table}[!htbp]
\centering
\footnotesize
\renewcommand{\arraystretch}{1.08}
\caption{Quadro comparativo do estado da arte e lacunas frente à proposta de mestrado.}
\begin{tabularx}{\textwidth}{p{3.0cm} p{2.9cm} X p{3.4cm}}
\toprule
\textbf{Referência} & \textbf{Domínio Técnico} & \textbf{Contribuição Principal} & \textbf{Lacunas e Limitações Identificadas} \\
Agüera y Arcas et al. / Google (2025) & Arquitetura e infraestrutura espacial de IA & Analisou a viabilidade e escala de data centers espaciais de IA alimentados por energia solar orbital. & Análise focada em escala de sistemas e economia macro, sem modelagem de controle térmico-elétrico dinâmico acoplado. \\
\addlinespace
Chen et al. / HotCarbon (2026) & Sistemas computacionais sustentáveis & Introduziu escalonamento térmico consciente de diafonia (\textit{thermal crosstalk}) para clusters orbitais de IA. & Tratou o subsistema elétrico como barramento ideal; omitiu perdas em conversores DC-DC e MPPT. \\
\addlinespace
Singh / arXiv (2026) & Alocação de cargas em nuvem orbital & Estruturou taxonomia e abstração semântica de cargas de IA viáveis para processamento orbital versus downlink. & Não modelou as dinâmicas físicas de transição térmica em vácuo nem a degradação de baterias. \\
\addlinespace
Pignaton et al. / ProPECAut (2021) & Modelagem de sistemas de energia solar & Validou modelo de 1 diodo / 5 parâmetros para módulos fotovoltaicos sob dados meteorológicos. & Desenvolvido exclusivamente para o espectro terrestre AM1.5 e resfriamento atmosférico natural. \\
\addlinespace
Souza, Silva e Almeida / CILAMCE (2021) & Eficiência energética em servidores (ProPECAut) & Aplicou métricas DCIM para otimizar consumo e climatização de data centers de médio porte. & Concepção estritamente terrestre baseada em condicionamento de ar convencional por convecção. \\
\addlinespace
Feng et al. / IEEE Access (2022) & Sistemas de armazenamento eletroquímico & Desenvolveu estimador de SoC via UKF com fator de esquecimento em células de lítio. & Validado apenas em bancada experimental, sem integração com despacho térmico de chips espaciais. \\
\addlinespace
Barros e Fernandes / REMUNOM (2026b) & Engenharia aeroespacial e sistemas orbitais & Estudo comparativo multifísico analítico entre LEO, SSO e GEO em 4 eixos com AHP. & Foco na seleção orbital macro e balanço anual; sem acoplamento de malha rápida de MPPT e transitórios kA/$\mu$s. \\
\addlinespace
Barros e Fernandes / REMUNOM (2026a) & Controle ótimo e inteligência artificial & Co-Design 4D em malha fechada via NMPC de 16 estados unindo térmico, elétrico, ADCS e downlink. & Validado para nós individuais de borda; sem detalhamento de redes de distribuição de potência em escala MW. \\
\midrule
\textbf{Proposta Desta Dissertação (2026)} & \textbf{Engenharia de Controle e Automação} & \textbf{Framework computacional integrado unindo geração AM0, perdas em conversores, Stefan-Boltzmann e TEMS/AUKF.} & \textbf{Validação computacional \textit{in silico} (sem TVAC); nós térmicos concentrados e foco termo-elétrico (sem física atômica de TID/SEB).} \\
\bottomrule
\end{tabularx}
\label{tab:estado_arte}
\end{table}

A síntese crítica do Quadro~\ref{tab:estado_arte} ratifica a oportunidade científica e a originalidade da dissertação, ao mesmo tempo em que delimita com clareza as fronteiras metodológicas assumidas nesta pesquisa de mestrado. Reconhece-se, a priori, que a validação fundamenta-se na simulação numérica multifísica de alta fidelidade calibrada por dados consolidados de literatura e parâmetros de fabricantes COTS espaciais, postergando ensaios físicos em câmara de vácuo térmico (\textit{Thermal Vacuum Chamber} -- TVAC) ou validações em voo orbital para futuros desdobramentos de doutorado ou projetos em parceria com a indústria aeroespacial. Ademais, a modelagem térmica adota uma abordagem de parâmetros concentrados (\textit{lumped thermal model}) representativa do comportamento dinâmico macroscópico dos chips e radiadores, focando estritamente no balanço térmico-energético e no controle ótimo do barramento, sem adentrar na física microscópica da degradação atômica por dose ionizante total (\textit{Total Ionizing Dose} -- TID) ou eventos de destruição por radiação única (\textit{Single Event Burnout} -- SEB).
